\section{Kết luận}
\label{sec:conclusion}

Đồ án đã hoàn thành trọn vẹn mục tiêu xây dựng và đánh giá hệ thống hỏi đáp tin tức \textbf{NewsLens} dựa trên kiến trúc Sinh có Tăng cường Truy xuất (RAG). Đối diện với thách thức lớn của bài toán hỏi đáp trên ngữ liệu tin tức quy mô lớn (11.064 bài báo CNN và 1.152 câu hỏi bổ sung ngữ cảnh tìm kiếm thuộc bộ dữ liệu NewsQA) xuất phát từ hai nguồn sai sót độc lập giữa khâu truy xuất thông tin và khâu sinh câu trả lời, nhóm tác giả đã tiếp cận bài toán một cách có hệ thống cả về mặt kỹ thuật công nghệ lẫn phương pháp luận nghiên cứu thực nghiệm.

\textbf{Về mặt kỹ thuật và sản phẩm phần mềm:} Đồ án đã hiện thực hóa hoàn chỉnh ứng dụng NewsLens với giao diện người dùng React (màn hình chat hiển thị trích dẫn nguồn và bàn thử nghiệm truy xuất), backend FastAPI và thư viện dùng chung \texttt{newsqa\_rag}. Hệ thống cung cấp nền tảng xử lý dữ liệu NewsQA, lập chỉ mục kép (Chroma vector database và BM25S sparse index), truy xuất kết hợp tái xếp hạng và kết nối linh hoạt với các mô hình ngôn ngữ lớn.

\textbf{Về mặt kết quả thực nghiệm:} Thông qua quy trình giải đấu đa tầng được kiểm soát khắt khe, đồ án đã xác lập cấu hình tối ưu toàn diện cho hệ thống:
\begin{itemize}
  \item Ở tầng truy xuất, BGE-M3 Sparse kết hợp bộ tái xếp hạng \texttt{bge-reranker-large} trên các đoạn văn bản 512 token (độ chồng lặp 64 token) chứng minh ưu thế vượt trội so với dense retriever truyền thống, mang lại mức tăng $+0{,}1634$ nDCG@5 và đẩy tỷ lệ tìm thấy bằng chứng Hit@5 lên $95{,}73\%$ trên tập phát triển. Phân tích dữ liệu chuyên sâu (EDA) giải mã cơ chế bản chất: trong miền tin tức phụ thuộc nhiều vào các thực thể định danh rời rạc, cơ chế khớp từ chính xác của mô hình thưa khai thác hiệu quả các điểm neo từ hiếm (IDF $\ge 6{,}0$, xuất hiện ở $59{,}5\%$ câu hỏi) để thu hẹp không gian từ 22.766 đoạn xuống trung vị 25 đoạn cạnh tranh mà không bị hiện tượng pha loãng thực thể (entity dilution) của dense bi-encoder; và bộ tái xếp hạng cross-encoder đóng vai trò quyết định để phân xử chính xác nhóm ứng viên cạnh tranh còn lại này.
  \item Ở tầng sinh đáp án, Prompt P2 (ràng buộc câu trả lời vào đúng loại thông tin trong một câu ngắn duy nhất) kết hợp 5 đoạn ngữ cảnh trên mô hình Gemini 3.1 Flash-Lite giúp tăng Answer Correctness lên $0{,}8015$ (tăng $+0{,}1043$ so với đối chứng), đồng thời bảo toàn hoàn hảo tính trung thực (Faithfulness $0{,}9603$) và tính hợp lệ của trích dẫn nguồn (Citation Validity $0{,}9786$).
  \item Kiểm định độc lập trên 284 câu hỏi của 50 bài báo held-out hoàn toàn chưa từng dùng để tinh chỉnh xác nhận hệ thống đạt Answer Correctness $0{,}7157$, chứng minh năng lực tổng quát hóa vững chắc trong môi trường dữ liệu thực tế.
\end{itemize}

\textbf{Về mặt phương pháp luận nghiên cứu:} Đồ án làm rõ giá trị của quy trình \textbf{đăng ký trước rào chắn an toàn (pre-registered guardrails)} và phương pháp kiểm định thống kê bootstrap ghép cặp gom cụm theo bài báo. Hiện tượng ``Luật thắng điểm số'' diễn ra qua ba giai đoạn thử nghiệm cho thấy: việc lựa chọn cấu hình thuần túy dựa vào điểm số cao nhất của một chỉ số đơn lẻ dễ dẫn đến việc chấp nhận các cấu hình có chất lượng không ổn định hoặc làm suy giảm các tiêu chí quan trọng khác (như độ trung thực, độ hợp lệ trích dẫn hay cấu trúc câu trả lời). Việc áp dụng hệ thống rào chắn an toàn đa chỉ số là bài học then chốt để đảm bảo hệ thống RAG đạt hiệu năng cân bằng và đáng tin cậy.

Hệ thống NewsLens không chỉ là một bài tập lớn môn học khai phá văn bản mà còn đặt nền móng vững chắc cho các ứng dụng hỏi đáp tin tức, thẩm định tin giả và tra cứu tư liệu báo chí trong thực tiễn.
